\section*{Chapitre \ref{ch:rs:anatomy-of-a-predictor} --- Anatomie d'un prédicteur}

\begin{solution}{exo:rs:anatomy-of-a-predictor:1}
$1/\sqrt{799} = 0.0354$. Une moyenne de 0.01 atteint $t = 2$ lorsque $0.01\sqrt n/0.0354 = 2$, soit $n = 50$ dates indépendantes, si le seul bruit
était l'échantillonnage transversal~; les chocs communs rendent les séries réelles d'IC plus bruitées, et il faut en pratique plusieurs fois autant de dates.
\end{solution}

\begin{solution}{exo:rs:anatomy-of-a-predictor:2}
Toujours 0.02 avec le rendement en excès de marché~: retrancher un même nombre à chaque rendement laisse les rangs inchangés. Face
aux rendements ajustés du bêta, la valeur change dès que les bêtas diffèrent, car chaque action retranche une quantité différente~; l'ampleur de la variation dépend de la façon dont
le \omterm{def:rs:anatomy-of-a-predictor:predictor}{prédicteur} est corrélé au bêta.
\end{solution}

\begin{solution}{exo:rs:anatomy-of-a-predictor:3}
$0.012/0.08 \times \sqrt{252} = 2.38$.
\end{solution}

\begin{solution}{exo:rs:anatomy-of-a-predictor:4}
Le résidu conserve $\rho = 0.6$ de la variance~: $0.03/\sqrt{0.6} = 0.039$.
\end{solution}

\begin{solution}{exo:rs:anatomy-of-a-predictor:5}
Jusqu'à $\sqrt{20} = 4.47$~; une valeur naïve de 6.0 peut n'être que $6.0/4.47 = 1.34$.
\end{solution}

\begin{solution}{exo:rs:anatomy-of-a-predictor:6}
Le \omterm{def:rs:anatomy-of-a-predictor:ic}{coefficient d'information de rang} (rank IC)~: robuste aux valeurs aberrantes du \omterm{def:rs:anatomy-of-a-predictor:predictor}{prédicteur} et des rendements, si bien qu'un seul titre extrême ne décide pas d'une date. Le \omterm{def:rs:anatomy-of-a-predictor:ic}{coefficient d'information} de Pearson~: un portefeuille dont
les poids sont proportionnels au \omterm{def:rs:anatomy-of-a-predictor:predictor}{prédicteur} gagne en proportion de la corrélation de Pearson, ampleurs comprises, si bien qu'il mesure ce que
produirait un portefeuille linéaire.
\end{solution}

\begin{solution}{exo:rs:anatomy-of-a-predictor:7}
Sur le dixième des jours où le rendement absolu de marché est le plus fort, le \omterm{def:rs:anatomy-of-a-predictor:ic}{coefficient d'information de rang} vaut 0.037 ($t = 5.0$)~; sur les autres, 0.039 ($t = 17.1$). Le
retournement est réparti également~: le simulateur inverse une fraction fixe du choc spécifique de chaque jour, quel que soit le comportement du marché. Sur des données réelles
il s'agit d'une dépendance à tester, non à supposer.
\end{solution}

\begin{solution}{exo:rs:anatomy-of-a-predictor:8}
Les cibles à 60 jours de deux jours consécutifs se recouvrent sur 59 jours, si bien que les 1\,260 IC sont loin d'être indépendants~: il y a environ 21
périodes non chevauchantes. La proposition borne la surestimation par $\sqrt{60}$~: le $t$ honnête peut n'être que $14/7.75 = 1.8$. Utiliser
Newey--West avec 59 décalages, ou retenir une date sur 60.
\end{solution}

\begin{solution}{pb:rs:anatomy-of-a-predictor:1}
\textbf{1.} 0.0385, $t = 17.8$. \textbf{2.} 0.0074, $t = 2.4$. \textbf{3.} 0.0162 (un jour contre cinq) et 0.0187 (cinq jours contre
un jour). \textbf{4.} Un jour contre un jour~: le simulateur inverse une fraction du choc spécifique de chaque jour le lendemain, et chacune des autres
paire ajoute quatre jours de rendements sans rapport avec le \omterm{def:rs:anatomy-of-a-predictor:predictor}{prédicteur} ou la cible. \textbf{5.} 0.0008~: rien. \textbf{6.} Les rangs ne changent pas
lorsqu'on retranche un même nombre à chaque rendement. \textbf{7.} 10~\% (le résidu conserve 90~\%). \textbf{8.} 52~\%~; $1/\sqrt{0.476} = 1.45$.
\textbf{9.} 0.0043 ($t = 0.8$), 0.0067 ($t = 3.5$), 0.0086 ($t = 3.5$). \textbf{10.} La composante sectorielle propre au \omterm{def:rs:anatomy-of-a-predictor:predictor}{prédicteur} ne porte aucune
information sur l'avenir dans ce marché, si bien que la retirer ôte du bruit au \omterm{def:rs:anatomy-of-a-predictor:predictor}{prédicteur} comme le résidu ôte du bruit
à la cible. \textbf{11.} 3.85 et 2.39. \textbf{12.} La borne $\sqrt h$ tient pour un \omterm{def:rs:anatomy-of-a-predictor:predictor}{prédicteur} persistant~; le retournement change chaque jour, si bien que
les IC de dates consécutives se recouvrent moins. \textbf{13.} $t = 1.31$~: valide, mais issu d'un cinquième des dates. \textbf{14.} $1/\sqrt{999} =
0.032$. \textbf{15.} 0.040, 6.4, 20.1. \textbf{16.} Pour un délai inférieur à la journée~: l'IC vaut 0.0385 avec le jour $t+1$ et 0.0008 avec le jour $t+2$.
\textbf{17.} Les modes de défaillance, puis l'IC lui-même~: davantage de fournisseurs de liquidité font disparaître le retournement par la concurrence (encombrement (crowding), chapitre~28).
\textbf{18.} \emph{Résultat de référence~:} pour un retournement à cinq jours évalué quotidiennement contre des cibles à cinq jours qui se recouvrent, le $t$ naïf vaut 3.85 et le
HAC $t$ 2.39~; sur le marché à forte composante sectorielle, neutraliser la cible et le \omterm{def:rs:anatomy-of-a-predictor:predictor}{prédicteur} fait passer l'IC de 0.0043 à 0.0086 et son $t$
de 0.8 à 3.5. \textbf{19.} La formule et la fenêtre du \omterm{def:rs:anatomy-of-a-predictor:predictor}{prédicteur}, la cible et son horizon, la normalisation et la \omterm{def:rs:anatomy-of-a-predictor:neutralise}{neutralisation},
l'univers et la période, le type d'IC et l'erreur type utilisée. \textbf{20.} Une formule avec une cible, un horizon, une normalisation et
une \omterm{def:rs:anatomy-of-a-predictor:neutralise}{neutralisation}, dont la formule est la partie la moins informative.
\end{solution}

\begin{solution}{iq:rs:anatomy-of-a-predictor:1}
La corrélation transversale entre un \omterm{def:rs:anatomy-of-a-predictor:predictor}{prédicteur} et le rendement qu'il prévoit, calculée date par date puis moyennée. Pour les rendements d'actions, un
IC moyen de 0.02 à 0.05 est bon s'il est stable~; le \omterm{def:rs:anatomy-of-a-predictor:icir}{ratio d'information de l'IC} et sa persistance comptent autant que sa
moyenne.
\iqlookfor{une définition par date, des ordres de grandeur réalistes et une attention à la stabilité.}
\end{solution}

\begin{solution}{iq:rs:anatomy-of-a-predictor:2}
Pas en l'état~: les cibles à 20 jours de dates consécutives se recouvrent sur 19 jours, si bien que le $t$ naïf peut surestimer d'un facteur allant jusqu'à $\sqrt{20} = 4.5$~; $t = 8$ pourrait n'être que
environ 1.8. Recalculer avec Newey--West (19 décalages) ou retenir une date sur 20.
\iqlookfor{la correction du recouvrement, chiffrée.}
\end{solution}

\begin{solution}{iq:rs:anatomy-of-a-predictor:3}
Pour retirer une composante qui ne porte aucune information mais ajoute du bruit (mouvements sectoriels), et éviter un portefeuille qui est en
réalité un pari sectoriel. Non lorsque l'information du \omterm{def:rs:anatomy-of-a-predictor:predictor}{prédicteur} se situe au niveau sectoriel (momentum sectoriel), ou lorsque le portefeuille doit prendre un
risque sectoriel.
\iqlookfor{la \omterm{def:rs:anatomy-of-a-predictor:neutralise}{neutralisation} comme affirmation sur l'endroit où se trouve l'information.}
\end{solution}

\begin{solution}{iq:rs:anatomy-of-a-predictor:4}
L'essentiel de sa puissance apparente vient de sa corrélation avec les facteurs retirés (sectoriel, bêta, taille)~: c'est en partie un pari sur les facteurs. Le 0.02 représente son habileté propre aux titres~; le reste en vaut la peine
dans la mesure où dépend de savoir si cette exposition aux facteurs est
rémunérée et recherchée.
\iqlookfor{la décomposition en parties factorielle et spécifique.}
\end{solution}

\begin{solution}{iq:rs:anatomy-of-a-predictor:5}
Pearson, si les positions doivent être proportionnelles à la prévision (les ampleurs comptent), en contrôlant en parallèle le rank IC pour la robustesse~; le rang
si la prévision ne sert qu'à ordonner les titres ou si ses queues ne sont pas fiables.
\iqlookfor{le lien entre la mesure et la façon dont la prévision est utilisée.}
\end{solution}

\begin{solution}{iq:rs:anatomy-of-a-predictor:6}
En l'absence de prévisibilité et avec des titres indépendants, la corrélation empirique de $N$ paires a une variance d'environ $1/(N - 1)$. La moyenne de $n$
dates indépendantes a pour erreur type $1/\sqrt{n(N - 1)}$~; $t = 2$ pour une moyenne de 0.02 avec $N = 1\,000$ exige $n = (2/(0.02 \times
31.6))^2 = 10$ dates en théorie, mais les chocs communs rendent la dispersion véritable de l'IC plusieurs fois plus grande, et, en pratique, il faut des centaines de
dates.
\iqlookfor{le résultat en $1/\sqrt{N-1}$ et la conscience qu'il sous-estime le bruit réel.}
\end{solution}
